\documentclass{article}
\usepackage{tcee-repaso}
\tcetema[El SEI hasta 1973]{3.B.20}{El sistema económico internacional desde el siglo XIX hasta la ruptura del sistema de Bretton-Woods.}
\tcefecha{13/11/2026}
\begin{document}
\cabecerarepaso
\begin{columnas}

% ==========================================================================
\bloque{Esquema del cante}
\esquemacante{%
\apartadocante{Introducción: «no hay bien ni mal que cien años dure»}{2,5}
\apartadocante{I. El SEI hasta 1914}{8}
\subcante{1. Sistema comercial}{2}
\subcante{2. Sistema financiero}{1,5}
\subcante{3. Sistema monetario: el patrón oro}{4,5}
\apartadocante{II. El SEI de entreguerras}{8}
\subcante{1. La Primera Guerra Mundial (1914-1929)}{3,5}
\subcante{2. La Gran Depresión (1929-1939)}{4,5}
\apartadocante{III. El sistema de Bretton Woods}{10,5}
\subcante{1. Los acuerdos y su funcionamiento}{6,5}
\subcante{\hspace*{0.8em}3 pilares y 3 planos}{2,5}
\subcante{\hspace*{0.8em}Ajuste, liquidez y confianza}{3}
\subcante{\hspace*{0.8em}Factores que lo sostuvieron}{1}
\subcante{2. La ruptura (1971-1973)}{4}
\apartadocante{Conclusión}{1}}
Funciones clave: paridad de acuñación $E_0=P^{\$}_{\mathrm{oro}}/P^{\pounds}_{\mathrm{oro}}=4{,}8665$; puntos del oro $E_0(1-c)\leq E\leq E_0(1+c)$; \autor{Triffin}: $D_t>G_t$. Gráficos: cronología, comercio mundial, puntos del oro, dilema de \autor{Triffin}.

% ==========================================================================
\bloque[2,5]{Introducción}
\begin{itemize}
\item \textbf{Enganche}: «no hay bien ni mal que cien años dure»: patrón oro, unos 40 años; cambios-oro, 6 (1925-1931); Bretton Woods, 27 (13 con convertibilidad europea). 15/8/1971: \autor{Nixon} cierra la ventanilla del oro. ¿Por qué se rompe cada orden y qué aprende el siguiente?
\item \textbf{Relevancia}: GATT/OMC, FMI, Banco Mundial y el papel del dólar nacen aquí; la historia, laboratorio (\autor{Schumpeter}, \autor{Samuelson}; Nobel 1993 \autor{North} y \autor{Fogel}; 2025 \autor{Mokyr}); «esta vez es distinto» (\autor{Reinhart}-\autor{Rogoff}); \autor{Santayana}.
\item \textbf{Contextualización}: \concepto{SEI}: instituciones, normas y prácticas de las relaciones económicas entre países. 3 \concepto{planos}: comercial, financiero, monetario. 3 \concepto{problemas} de un sistema monetario: ajuste, liquidez, confianza.
\item \textbf{Problemática}: ¿cómo funcionaba el oro y por qué cayó? ¿Qué se pactó en 1944 y por qué no aguantó? ¿Qué tienen en común las caídas?
\item \textbf{Estructura}: orden histórico, 3 planos en cada etapa: \hyperref[rep:I]{I}, \hyperref[rep:II]{II}, \hyperref[rep:III]{III}.
\end{itemize}
\graficorepaso[\linewidth]{cronologia-sei}{1815-1973: regímenes monetarios en el centro, comercio arriba, dinero abajo}

% ==========================================================================
\bloque[8]{I. El SEI hasta 1914\label{rep:I}}
\begin{ideaclave}
Primera globalización sin tratado (cooperación \emph{de facto}): librecambio, exportación de capital y personas y patrón oro, que funcionó por credibilidad, cooperación y hegemonía británica más que por su mecanismo automático.
\end{ideaclave}
\graficorepaso[0.95\linewidth]{comercio-mundial}{Comercio mundial: \dato{12\,\% (1870) $\to$ 21\,\% (1913) $\to$ 8\,\% (1932) $\to$ 30\,\% (1973) del PIB}{Klasing-Milionis; PWT}}

\sub{I.1. Sistema comercial\label{rep:I.1}}
\begin{itemize}
\item Ideas: \autor{Smith} (1776), \autor{Ricardo} (1817), \autor{Mill} (1848); \autor{Malthus}, por las Leyes de Granos.
\item 3 hitos: Tratado de Edén (1786); derogación de las Leyes de Granos (\autor{Robert Peel}, 1846; Navegación, 1849); \autor{Cobden}-\autor{Chevalier} (1860) con \concepto{nación más favorecida} incondicional $\Rightarrow$ red multilateral.
\item Vuelta al proteccionismo desde 1879: \autor{Bismarck} (1879), \autor{McKinley} (1890), \autor{Méline} (1892); cereal americano y depresión 1873-1896. RU librecambista hasta 1914.
\item Convergencia de precios y salarios (\autor{O'Rourke}-\autor{Williamson}; \autor{Heckscher}-\autor{Ohlin}).
\end{itemize}

\sub{I.2. Sistema financiero\label{rep:I.2}}
\begin{itemize}
\item RU con superávit $\Rightarrow$ Londres presta: bonos de Estados y ferrocarriles (cartera); activos exteriores $\approx$ 1/5 del PIB mundial (\autor{Obstfeld}-\autor{Taylor}).
\item \dato{60 millones de europeos emigran (1820-1914)}{Hatton-Williamson}. Préstamos de oro entre bancos centrales (Baring, 1890).
\end{itemize}

\sub{I.3. El patrón oro (a-e)\label{rep:I.3}}
\begin{itemize}
\item \emph{Idea}: del bimetalismo (Francia 15,5:1) y la ley de \autor{Gresham} al oro: RU \emph{de facto} 1717 (\autor{Newton}), ley 1816, convertible 1821; Alemania 1871-1873, EE.\,UU. 1879/1900.
\item \emph{Supuestos} (reglas del juego): paridad fija, convertibilidad plena, libre circulación del oro, no esterilizar, tipo de descuento.
\end{itemize}
\begin{formula}
\[E_0=\frac{20{,}67\ \$}{4{,}2477\ \pounds}=4{,}8665\ \$/\pounds\]\[E_0(1-c)\leq E\leq E_0(1+c)\]
\end{formula}
\graficorepaso[0.95\linewidth]{puntos-oro}{En $\overline E$ la oferta de libras es horizontal: sale oro de EE.\,UU.}
\begin{itemize}
\item \emph{Desarrollo}: flujo-especie (\autor{Hume}, 1752): déficit $\to$ sale oro $\to\downarrow M\to\downarrow P\to$ competitividad \vertema{3.B.12}. Precios: deflación 1873-1896, inflación 1896-1914.
\item \emph{Implicaciones}: fijo + capital libre = sin política monetaria propia (trilema). Ventajas: estabilidad, credibilidad (regla con cláusula de escape, \autor{Bordo}-\autor{Kydland}). Inconvenientes: subordinación, productores de oro, sesgo deflacionista.
\item \emph{Extensiones}: cambios-oro en la periferia; Banco de Inglaterra, «director de orquesta» (\autor{Keynes}).
\item \emph{Evidencia}: \dato{60\,\% de esterilización}{Bloomfield, 1959}; ajuste por capital y renta, no por oro y precios; periferia que suspende.
\item \emph{Valoración}: credibilidad y cooperación; ajuste ni automático ni simétrico.
\end{itemize}

% ==========================================================================
\bloque[8]{II. El SEI de entreguerras\label{rep:II}}
\begin{ideaclave}
La guerra rompe hegemonía, cooperación y prioridad del equilibrio externo. Volver al oro sin ellas fracasa y transmite la Gran Depresión; en los 30, cada uno se salva solo.
\end{ideaclave}

\sub{II.1. La Primera Guerra Mundial (1914-1929)\label{rep:II.1}}
\begin{itemize}
\item EE.\,UU. acreedor y «felices 20»; Europa, paro de reconversión.
\item Vencidos: inflación, escasez, \concepto{reparaciones} (Versalles, 1919; \dato{132.000 M de marcos oro}{1921}; \autor{Keynes}, 1919). Ruhr (1923) $\to$ \dato{4,2 billones de marcos por dólar}{11/1923}; \emph{Rentenmark}.
\item Plan \autor{Dawes} (1924; préstamo de 800 M de marcos oro), Plan \autor{Young} (1929; BPI, 1930), moratoria \autor{Hoover} (1931), Lausana (1932) $\to$ Plan \autor{Brady} \vertema{3.B.29}.
\item Comercio: proteccionismo, bloques, desequilibrios financiados con capital de EE.\,UU.
\item Monetario: Génova (1922) $\to$ \concepto{patrón de cambios-oro}; EE.\,UU. 1919, RU 1925 a 4,86 (lingotes; \autor{Keynes} contra \autor{Churchill}), Francia 1926/1928 infravalorada. Falla: sin centro, sin cooperación (EE.\,UU. y Francia esterilizan oro), sin credibilidad (sufragio y sindicatos), liquidez mal repartida.
\end{itemize}

\sub{II.2. La Gran Depresión (1929-1939)\label{rep:II.2}}
\begin{itemize}
\item EE.\,UU. 1929-1933: \dato{PIB $-26$\,\%, precios $-25$\,\%, paro 25\,\%}{BEA, BLS}; comercio mundial a un tercio (espiral de \autor{Kindleberger}).
\item Causas: sobreproducción o subconsumo (\autor{Keynes}); dinero $-35$\,\% (\autor{Friedman}-\autor{Schwartz}, 1963; Nobel 1976); quiebras y crédito (\autor{Bernanke}, 1983; Nobel 2022); deflación por deudas (\autor{Fisher}, 1933); austeridad. Detonante: crac de octubre de 1929.
\item \textbf{Propagación}: el oro como correa de transmisión (\emph{Golden Fetters}, \autor{Eichengreen}, 1992); Credit-Anstalt (5/1931); libra fuera el 21/9/1931: crisis de \emph{segunda generación} (\autor{Eichengreen}-\autor{Jeanne}).
\item Salida: quien deja antes el oro se recupera antes (\autor{Eichengreen}-\autor{Sachs}, 1985): RU 1931, EE.\,UU. 1933 (35~\$/onza, 1934), bloque oro 1936. \emph{New Deal} (3 R).
\item Consecuencias: \autor{Smoot}-\autor{Hawley} (\dato{59\,\% en 1932}{Irwin}), Ottawa (1932), RTAA (1934); impagos, controles, Glass-Steagall (1933); devaluaciones competitivas, bloques, Acuerdo Tripartito (1936). \autor{Kindleberger}: faltó un estabilizador.
\end{itemize}

% ==========================================================================
\bloque[10,5]{III. El sistema de Bretton Woods\label{rep:III}}
\begin{ideaclave}
Primer orden diseñado, con instituciones en los 3 planos: \dato{PIB mundial +4,9\,\% al año 1950-1973}{Maddison}. El plano monetario nace con los 3 problemas y descansa en el dólar; cae cuando el capital se mueve y EE.\,UU. antepone sus objetivos internos.
\end{ideaclave}

\sub{III.1. Los acuerdos y su funcionamiento\label{rep:III.1}}
\begin{itemize}
\item Julio de 1944, 44 países. Plan \autor{White} (fondo de 5.000 M, ajuste del deudor) frente a Plan \autor{Keynes} (Unión de Compensación, \emph{bancor}, ajuste simétrico): gana \autor{White}; queda la moneda escasa.
\item \textbf{3 pilares}: FMI; BIRF; OIC fallida (La Habana, 1948) $\to$ GATT (1947, 23 países; NMF, trato nacional, transparencia, sin cuotas; Ronda Kennedy $-35$\,\%).
\item Financiero: Plan \autor{Marshall} (\dato{13.000 M\$, 90\,\% donaciones}{1948-1952}), OECE, UEP (1950-1958); controles de capital.
\item Monetario: paridades $\pm1$\,\%, dólar a 35~\$/onza (solo bancos centrales), «desequilibrio fundamental» (10\,\% acumulado), convertibilidad corriente (Europa, 12/1958); FMI con cuotas (25\,\% en oro). $n-1$ tipos. España: FMI 1958, 60 pta/\$ (1959), 70 (1967) \vertema{4.A.27}.
\item \textbf{Ajuste}: 3 asimetrías: deudor/acreedor; $n-1$ (EE.\,UU. sin trilema); señoreaje (\concepto{privilegio exorbitante}, \autor{Giscard d'Estaing}, 1965).
\item \textbf{\autor{Triffin}} (1960, a-e): liquidez solo con déficit de EE.\,UU. $\Rightarrow D_t\uparrow$, $G_t$ fijo.
\end{itemize}
\begin{formula}
\[D_t>G_t\ \Rightarrow\ \text{la convertibilidad deja de ser creíble}\]
\end{formula}
\graficorepaso[0.95\linewidth]{triffin}{\dato{Oro 24,6 (1949) $\to$ 10,2 (1971); $D_t>G_t$ en 1961; 66,5 en 1971}{Fed, Z.1; miles de M\$}}
\begin{itemize}
\item Liquidez o confianza; solución: activo supranacional (DEG). Ext.: banco del mundo (\autor{Kindleberger}-\autor{Despres}-\autor{Salant}); versión moderna (\autor{Farhi}-\autor{Gourinchas}-\autor{Rey}). Valoración: también la política expansiva (Vietnam).
\item \textbf{Confianza}: libra $-14{,}3$\,\% (1967); marco +5\,\% (1961) y +9,3\,\% (1969); franco $-11{,}1$\,\% (1969); ataques de primera generación (\emph{pool} 1968, 1971; \autor{Garber}).
\item Parches: \emph{pool} del oro (1961-1968), AGP (1962, 6.000 M), \emph{swaps}, compromiso alemán (1967), controles de EE.\,UU. (1963, 1965, 1968), DEG (1969/1970). \autor{Friedman} (1953) pide flotar.
\end{itemize}

\sub{III.2. La ruptura (1971-1973)\label{rep:III.2}}
\begin{itemize}
\item 15/8/1971 (\autor{Nixon}): fin de la convertibilidad, recargo del 10\,\%, congelación de precios y salarios.
\item Smithsonian (18/12/1971): 38~\$/onza ($-7{,}9$\,\%), bandas $\pm2{,}25$\,\% (serpiente, 1972), sin recargo. Febrero de 1973: 42,22~\$ ($-10$\,\%); marzo de 1973: flotación. Petróleo (10/1973). Jamaica (1976) \vertema{3.B.21}.
\item Cayó por ajuste, liquidez, confianza y hegemonía; balance: 25 años de crecimiento e instituciones vivas.
\end{itemize}

% ==========================================================================
\bloque[1]{Conclusión}
\begin{itemize}
\item \textbf{Recapitulación}: comercio (librecambio $\to$ crisis $\to$ GATT); finanzas (libertad $\to$ controles); dinero (oro $\to$ dólar), sin resolver a la vez ajuste, liquidez y confianza.
\item \textbf{Relevancia}: trilema, asimetría, privilegio del dólar, \autor{Triffin} moderno.
\item \textbf{Extensiones}: \vertema{3.B.21} \vertema{3.B.22} \vertema{3.B.17} \vertema{3.B.15} \vertema{3.B.16} \vertema{4.A.27}.
\item \textbf{Opinión}: caen cuando falta la cooperación y la voluntad política.
\item \textbf{Idea final}: «no hay bien ni mal que cien años dure»: dura el sistema que se adapta y en el que se coopera.
\end{itemize}

% ==========================================================================
\bloque{Los 3 órdenes por planos}
{\small\noindent\begin{tabular}{@{}>{\raggedright\arraybackslash}p{1.5cm}|>{\raggedright\arraybackslash}p{2cm}|>{\raggedright\arraybackslash}p{2cm}|>{\raggedright\arraybackslash}p{2cm}@{}}
\textbf{Plano} & \textbf{Hasta 1914} & \textbf{Entreguerras} & \textbf{Bretton Woods}\\ \hline
Comercial & librecambio, NMF; proteccionismo desde 1879 & aranceles, bloques & GATT y CEE\\ \hline
Financiero & capital y emigración libres & reparaciones, impagos, controles & BIRF, Marshall, controles\\ \hline
Monetario & oro (Londres) & cambios-oro, flotación & paridades al dólar, dólar al oro\\
\end{tabular}\par}

% ==========================================================================
\bloque{Definiciones}
\begin{itemize}
\item \concepto{Sistema económico internacional}: instituciones, normas y prácticas que ordenan comercio, finanzas y pagos entre países.
\item \concepto{Bimetalismo}: unidad definida en oro y plata con relación fija. \concepto{Ley de \autor{Gresham}}: la moneda mala (sobrevalorada) expulsa a la buena; compara paridad oficial y de mercado.
\item \concepto{Patrón oro}: paridad fija en oro y convertibilidad plena. \concepto{Patrón de cambios-oro}: reservas en oro y en divisas convertibles (Génova, 1922; Bretton Woods). \concepto{Patrón lingotes-oro}: solo se convierten lingotes (RU, 1925).
\item \concepto{Reglas del juego}: no esterilizar y mover el tipo de descuento con el oro. \concepto{Esterilizar}: compensar con activos internos los movimientos de reservas.
\item \concepto{Puntos del oro}: tipos a los que conviene arbitrar con oro (paridad $\pm$ coste de envío).
\item \concepto{Desequilibrio fundamental}: condición para cambiar la paridad en Bretton Woods; nunca se definió.
\item \concepto{Problema de los $n-1$}: con $n$ monedas, solo $n-1$ tipos; la moneda $n$ no interviene.
\item \concepto{Privilegio exorbitante}: señoreaje internacional del emisor de la moneda de reserva.
\item \concepto{DEG}: activo de reserva del FMI (1969); hoy, cesta de dólar, euro, renminbi, yen y libra.
\item \concepto{Cláusula de nación más favorecida}: toda ventaja concedida a uno se extiende a todos (art.~I del GATT).
\end{itemize}

% ==========================================================================
\bloque{Datos clave}
\begin{itemize}
\item Comercio/PIB mundial: \dato{12\,\% (1870), 21\,\% (1913), 8\,\% (1932), 20\,\% (1950), 30\,\% (1973)}{Klasing-Milionis; PWT}.
\item PIB mundial: \dato{+2,1\,\% al año (1870-1913); +4,9\,\% (1950-1973)}{Maddison, 2001}.
\item Emigración: \dato{60 millones de europeos, 1820-1914}{Hatton-Williamson}.
\item Gran Depresión en EE.\,UU. (1929-1933): \dato{PIB $-26$\,\%; dinero $-35$\,\%; paro 25\,\%; Dow Jones $-89$\,\% (1929-1932)}{BEA; F-S}. Importaciones de 75 países: \dato{2.998 $\to$ 992 M\$ oro al mes (1/1929 $\to$ 1/1933)}{Kindleberger}.
\item Reparaciones: \dato{132.000 M de marcos oro (1921) $\to$ $\approx$ 112.000 M (Young, 1929)}{Comisión de Reparaciones}.
\item FMI: \dato{cuotas iniciales de 8.800 M\$, con la URSS}{Convenio, 1944}. AGP: \dato{6.000 M\$}{FMI, 1962}.
\item Oro de EE.\,UU.: \dato{24,6 mm\$ (1949) $\to$ 10,2 (1971)}{Fed, Z.1}; pasivos líquidos frente al exterior: \dato{8,3 (1950) $\to$ 66,5 (1971)}{Fed, Z.1}; oro en Londres a \dato{$\approx$40~\$ (10/1960)}{mercado de Londres}.
\end{itemize}

% ==========================================================================
\bloque{Bretton Woods en 5 reglas}
\begin{enumerate}
\item Paridades fijas pero ajustables frente al dólar, bandas del $\pm1$\,\%.
\item Dólar convertible en oro a 35~\$/onza, solo para bancos centrales.
\item Cambios de paridad ante «desequilibrio fundamental»: hasta un 10\,\% sin objeción del FMI.
\item Convertibilidad por cuenta corriente (art.~VIII); controles de capital permitidos (art.~VI).
\item FMI: cuotas (25\,\% en oro), voto y préstamos condicionados; cláusula de moneda escasa.
\end{enumerate}

% ==========================================================================
\bloque{Autores y fechas\label{rep:autores}}
\hito{1717}{\autor{Newton}: el RU, de hecho, al oro}
\hito{1752}{\autor{Hume}: flujo-especie}
\hito{1846}{Derogación de las Leyes de Granos (\autor{Peel})}
\hito{1860}{\autor{Cobden}-\autor{Chevalier}}
\hito{1879}{EE.\,UU. al oro; arancel de \autor{Bismarck}}
\hito{1919}{Versalles; \autor{Keynes}: \emph{Consecuencias}}
\hito{1925}{La libra vuelve al oro (Génova, 1922)}
\hito{1930}{\autor{Smoot}-\autor{Hawley}; BPI}
\hito{1931}{Credit-Anstalt; la libra sale del oro}
\hito{1944}{Bretton Woods}
\hito{1947}{GATT; Plan \autor{Marshall} (1948)}
\hito{1953}{\autor{Friedman}: tipos flexibles}
\hito{1958}{Convertibilidad europea}
\hito{1959}{\autor{Bloomfield}; \autor{Triffin} en el Congreso}
\hito{1960}{\autor{Triffin}: \emph{Gold and the Dollar Crisis}}
\hito{1961}{\emph{Pool} del oro; AGP (1962)}
\hito{1963}{\autor{Friedman}-\autor{Schwartz}}
\hito{1968}{Doble mercado del oro; DEG (1969)}
\hito{1971}{\autor{Nixon}; Smithsonian}
\hito{1973}{Flotación; \autor{Kindleberger}}
\hito{1985}{\autor{Eichengreen}-\autor{Sachs}}
\hito{1992}{\autor{Eichengreen}: \emph{Golden Fetters}}

% ==========================================================================
\bloque{Las lecturas de la Gran Depresión}
{\small\noindent\begin{tabular}{@{}>{\raggedright\arraybackslash}p{2.3cm}|>{\raggedright\arraybackslash}p{5.6cm}@{}}
\textbf{Lectura} & \textbf{Qué falló}\\ \hline
Keynesiana & la demanda (inversión, consumo); salida fiscal \vertema{3.A.4}\\ \hline
Monetarista & la Fed dejó caer el dinero un 35\,\% (\autor{Friedman}-\autor{Schwartz}) \vertema{3.A.5}\\ \hline
Financiera & quiebras y crédito (\autor{Bernanke}, 1983); deflación por deudas (\autor{Fisher}) \vertema{3.B.26}\\ \hline
Internacional & el oro transmitió la contracción (\autor{Eichengreen}); faltó un estabilizador (\autor{Kindleberger})\\ \hline
Austriaca & sobreinversión previa: había que liquidar \vertema{3.A.3}\\
\end{tabular}\par}

% ==========================================================================
\bloque{Patrón oro clásico frente a Bretton Woods}
{\small\noindent\begin{tabular}{@{}>{\raggedright\arraybackslash}p{2cm}|>{\raggedright\arraybackslash}p{2.9cm}|>{\raggedright\arraybackslash}p{2.9cm}@{}}
 & \textbf{Patrón oro} & \textbf{Bretton Woods}\\ \hline
Ancla & oro (Londres, centro) & dólar, y el dólar al oro\\ \hline
Tipos & fijos, puntos del oro & fijos pero ajustables, $\pm1$\,\%\\ \hline
Capital & libre & controlado hasta los 60\\ \hline
Trilema & fijo + capital: sin política propia & fijo + política propia: con controles\\ \hline
Ajuste & oro, capital y renta; asimétrico & reservas, FMI, realineamientos tardíos\\ \hline
Liquidez & oro y libras & dólares (déficit de EE.\,UU.) y DEG\\ \hline
Fin & guerra de 1914 & 1971-1973\\
\end{tabular}\par}

% ==========================================================================
\bloque{Lo que cada orden aprendió del anterior}
\begin{itemize}
\item \emph{De 1914 a los años 20}: se quiso volver al oro y a la paridad de antes; se ahorró oro con divisas (Génova), pero sin acuerdo ni centro.
\item \emph{De los años 30 a 1944}: contra las devaluaciones competitivas, paridades vigiladas por el FMI; contra la falta de reservas, un fondo de cuotas; contra el proteccionismo, el GATT y la nación más favorecida; contra la rigidez del oro, paridades ajustables y controles de capital.
\item \emph{De 1973 en adelante}: flotación de las grandes monedas, el FMI como supervisor y la búsqueda de estabilidad regional en Europa (serpiente, SME, euro) \vertema{3.B.16} \vertema{3.B.21}.
\end{itemize}

% ==========================================================================
\bloque{Preguntas del tribunal}
\begin{enumerate}
\item \textbf{¿Por qué funcionó el patrón oro antes de 1914 y no después?} Credibilidad (sufragio restringido), cooperación y centro británico; tras 1918, nada de eso.
\item \textbf{¿El ajuste del patrón oro era por precios?} En teoría (\autor{Hume}); en la práctica, por capital y renta, con esterilización (\autor{Bloomfield}).
\item \textbf{¿Qué es un punto del oro?} El tipo al que conviene mandar oro en vez de comprar divisas: paridad $\pm$ coste de transporte.
\item \textbf{¿Qué causó la Gran Depresión?} Varias cosas; la dimensión internacional: el oro transmitió la contracción y salir de él permitió recuperarse.
\item \textbf{¿Por qué la crisis de la libra de 1931 es de segunda generación?} Quedaban reservas; el paro hacía creíble el abandono.
\item \textbf{¿Qué es el dilema de \autor{Triffin}?} Liquidez exige déficit de EE.\,UU.; el déficit mina la confianza en la convertibilidad.
\item \textbf{¿Por qué EE.\,UU. no podía devaluar?} Es la moneda $n$: solo si los demás revaluaban.
\item \textbf{¿Por qué se impuso el Plan \autor{White}?} EE.\,UU. era el acreedor; \autor{Keynes} representaba a un deudor.
\item \textbf{¿Por qué no funcionó el patrón de cambios-oro de los 20?} Paridades mal elegidas (libra sobrevalorada, franco infravalorado), esterilización del oro en EE.\,UU. y Francia, y reservas en libras que dependían de la confianza.
\item \textbf{¿Qué relación hay entre Bretton Woods y el trilema?} Eligió tipos fijos con política propia gracias a los controles; al moverse el capital, el lado se volvió insostenible \vertema{3.B.15}.
\item \textbf{¿Qué fue el \emph{pool} del oro?} 8 bancos centrales vendían oro en Londres para mantenerlo en 35~\$; se cerró en 1968 y llegó el doble mercado.
\item \textbf{¿Fue \autor{Triffin} o la inflación de EE.\,UU.?} Las 2: el dilema era estructural; Vietnam y la política expansiva lo precipitaron.
\item \textbf{¿Qué es la espiral de \autor{Kindleberger}?} El valor del comercio mundial mes a mes en 1929-1933: cada vuelta, más pequeña; cayó a un tercio.
\item \textbf{¿Por qué el Plan \autor{Marshall} fuera de Bretton Woods?} El BIRF era pequeño para reconstruir Europa; EE.\,UU. dio donaciones, con la OECE como gestora y la UEP para los pagos.
\item \textbf{¿Qué aportó la cláusula de nación más favorecida?} Multilateraliza las rebajas bilaterales: en 1860 (\autor{Cobden}-\autor{Chevalier}) y en el GATT (art.~I); a cambio, polizones.
\item \textbf{¿Qué sobrevivió a 1973?} Las instituciones (FMI, Banco Mundial, GATT/OMC), no las paridades.
\end{enumerate}

\begin{trampa}
\begin{itemize}
\item Leyes de Granos: las derogó \autor{Robert} Peel (no «Thomas»), 1846.
\item Paridad clásica: 4,8665~\$/\pounds{} (20,67~\$ y 4,2477~\pounds{} por onza fina), no «50~\pounds{} por onza».
\item Bandas: $\pm1$\,\% (1944-1971), $\pm2{,}25$\,\% (Smithsonian). Febrero de 1973: $-10$\,\% (42,22~\$).
\item Plan \autor{Marshall}: donaciones en su mayoría, no préstamos. Plan \autor{Dawes}: 800 M de \emph{marcos oro}.
\item \autor{Friedman}, \emph{The Case for Flexible Exchange Rates}: 1953 (no 1969).
\item Bancos en 1930-1933: más de una quinta parte suspendió pagos (no «un tercio quebró»).
\item DEG: Primera Enmienda (1969), primera asignación 1/1/1970; Jamaica (1976) y Segunda Enmienda (1978), ya en el 3.B.21.
\end{itemize}
\end{trampa}

\end{columnas}
\end{document}
